\documentclass[11pt,a4paper]{article}
\usepackage[T1]{fontenc}
\usepackage{lmodern}
\usepackage[a4paper,margin=2.5cm]{geometry}
\usepackage{microtype,amsmath,amssymb,graphicx,booktabs,tabularx,enumitem}
\usepackage[numbers]{natbib}
\usepackage[hidelinks]{hyperref}
\setlist{itemsep=1pt,topsep=3pt}

\title{Predicting risky choices for new participants}
\author{Kanchan Bharti \quad [5020261053] \quad kanchanbharti039\\
  \small Track B Behavioural Economics \quad
  Team: TDI Lake Grove\\
  \small CS-3410 / MAT-3211-1 Introduction to Machine Learning, Ashoka University}
\date{4 November 2026}

\begin{document}
\maketitle

\begin{abstract}
I ask whether the gain/loss framing of a gamble and expected-value (EV) features help to predict whether
a \emph{previously unseen} participant picks the risky option over a sure amount. The data are
\textbf{simulated} (8\,400 decisions from 210 participants). I compare an EV rule, four lecture models and
a logistic regression with engineered features under participant-grouped 5-fold cross-validation, and
test the effect of the split and of \texttt{response\_time\_ms}. EV features raise ROC-AUC from
$0.723\pm0.011$ (raw numbers) to $0.785\pm0.015$ and domain/frame to $0.814\pm0.014$; person traits add
nothing. The final rank-average of logistic regression and ExtraTrees reaches $0.820\pm0.011$
(fold mean $\pm$ s.d.), a gain of only $0.002$ over logistic regression alone. The main limitation is that the
data are simulated and noisy, so the estimates say nothing about real people and the last few thousandths
of AUC cannot be separated from fold noise.
\end{abstract}

\section{Research question}
\textbf{Question.} For new participants, do (i) the gain/loss frame and problem domain and (ii) features based on
expected value predict risky choice better than a model that only sees the raw problem numbers?
Predictions are scored by ROC-AUC on a held-out set of participants.

\textbf{Expectation.} I expected EV and frame features to improve AUC, and person details (age, numeracy,
self-reported risk taking) not to help for unseen people. The expectation would be \emph{contradicted} if the
frame/domain features added no AUC beyond EV features, or if person traits improved held-out AUC.
Two further expectations are tested in Section~\ref{sec:exps}: a random split is optimistic compared with a
participant split, and \texttt{response\_time\_ms} adds nothing.

\section{Background and related work}
Expected value is the classical benchmark for choice under risk, but people depart from it in systematic ways.
Prospect theory describes choices as depending on gains and losses relative to a reference point, with a value
function that is steeper for losses and a probability weighting that overweights small probabilities
\citep{kahneman1979,tversky1992}. The same outcomes presented as gains or as losses lead to different choices
\citep{tversky1981}, so the \emph{frame} should carry information that EV cannot. The simulator behind this
dataset is prospect-theory-like, so I build a few fixed prospect-theory value differences as features (no
parameter is fitted to data).

On the methodological side, observations from the same person are not independent, so the validation split must
follow the structure of the data; otherwise performance is overestimated \citep{roberts2017}. I report ROC-AUC,
which depends only on the ranking of scores \citep{fawcett2006}. Besides lecture models I use Extremely
Randomized Trees \citep{geurts2006}, via scikit-learn \citep{pedregosa2011}.

\section{Data and preprocessing}
\textbf{Data.} The Kaggle ``Track B: risky choice'' data (\emph{simulated}): \texttt{train.csv} has 8\,400
decisions from 210 participants (40 each) and \texttt{test.csv} 3\,600 decisions from 90 \emph{different}
participants (no overlap). Each row is one choice between a risky gamble (probability \texttt{prob\_high} of
\texttt{risky\_high\_inr}, otherwise \texttt{risky\_low\_inr}) and a sure amount. The target
\texttt{chose\_risky} is balanced (48.2\% risky). The four combinations of domain and frame are
gain/gain-frame, gain/loss-frame, loss/none and mixed/none; the frame manipulation exists only within the gain
domain.

\textbf{Audit and exclusions.} The only missing values are in \texttt{realised\_payout\_inr} (4\,600 of 8\,400),
which is the money paid \emph{after} the choice and is absent from the test file; it is dropped as leakage.
\texttt{response\_time\_ms} is measured while the person decides, so it is not available before the choice; it is
excluded from all final models and used only in a control experiment. \texttt{trial\_index} is known in advance
and is allowed (used by the tree model only).

\textbf{Features.} All formulas use only the problem numbers, so nothing is learned from the data:
the EV difference $\mathrm{EV}=p\,x_H+(1-p)\,x_L-s$ and quantities derived from it (EV divided by the
stakes spread and by the endowment, the position of the sure amount within the gamble, indicators for
$p\le0.1$ and $p\ge0.9$); four fixed prospect-theory value differences (parameter sets
$(\alpha,\lambda,\gamma)$ taken from the literature range, normalised by the endowment); and interactions of EV and
prospect-theory scores with the domain/frame cell, and of EV with person traits. This gives 48 columns for the
linear model and 24 for the tree model. Scaling and one-hot encoding are inside the scikit-learn pipeline and are
fitted on training folds only.

\section{Methods and experimental design}
\subsection{Validation protocol}
Test participants are all new, so I use \texttt{GroupKFold} on \texttt{participant\_id}: the 40 decisions of a
person never appear in both training and validation folds. I use $k=5$ shuffled folds with seed 0, and the
\emph{same} folds for every model, which allows paired comparisons. I report the mean and standard deviation
(over the 5 folds) of AUC, log-loss and accuracy at threshold 0.5. The only tuned quantity was the
L2 strength $C$ of logistic regression, scanned over $\{0.003,0.03,0.1,1\}$ on the same folds; AUC was flat
(Table~\ref{tab:appC}) and I fixed $C=0.1$. The final model choice was made on the same cross-validation, so the
estimate may be slightly optimistic. No test labels were used.

\subsection{Baselines}
(1) A constant prior predictor (trivial baseline: AUC must be 0.5). (2) The EV rule, scoring each decision by
the EV difference, with no training; this is the economic benchmark. (3)--(6) Logistic regression, LDA, Gaussian
Naive Bayes and a linear SVM from the lectures, each on the four raw problem columns only. They show what
standard models achieve without feature engineering.

\subsection{Controlled experiments}\label{sec:exps}
\begin{enumerate}
\item \textbf{Experiment 1, feature groups.} \emph{Change:} feature groups are added one at a time to a logistic
  regression: (a) raw numbers; (b) + EV features; (c) + domain/frame; (d) + person and session traits;
  (e) + frame$\times$EV interactions; (f) + prospect-theory and trait interactions. \emph{Control:} model class,
  folds and metric. \emph{Expectation:} (b) and (c) help; (d) does not.
\item \textbf{Experiment 2, split scheme.} \emph{Change:} random 5-fold split versus participant-grouped split,
  with the final logistic regression, and again after adding a one-hot \texttt{participant\_id}.
  \emph{Control:} model, features, number of folds. \emph{Expectation:} the random split is optimistic, especially
  when the model can identify the person.
\item \textbf{Experiment 3, response time.} \emph{Change:} add \texttt{response\_time\_ms} to the final
  logistic regression. \emph{Control:} everything else. \emph{Expectation:} no gain, so excluding it costs nothing.
\end{enumerate}
Finally I compare ExtraTrees (300 trees, minimum leaf 30, 60\% of features per split) and the \textbf{rank
average} of the final logistic regression and ExtraTrees. Because AUC depends only on ordering, averaging ranks is a
safe way to combine models; the submitted values are therefore scores in $[0,1]$, not calibrated probabilities.

\section{Results}
Table~\ref{tab:main} gives the main results; Figure~\ref{fig:overview} shows the framing effect, ROC curves and
calibration; paired differences are in Table~\ref{tab:paired} (Appendix).

\begin{table}[ht]
\centering
\caption{Participant-grouped 5-fold validation (seed 0). Mean $\pm$ standard deviation over 5 folds. A dash means the
model gives no probabilities (log-loss) or no threshold (accuracy). Rows in bold: final logistic regression and
submitted blend.}\label{tab:main}
\small
\begin{tabularx}{\linewidth}{@{}Xccc@{}}
\toprule
Model (features) & ROC-AUC & Log-loss & Accuracy \\
\midrule
Constant prior (trivial)            & $0.500\pm0.000$ & $0.693\pm0.002$ & $0.518\pm0.032$ \\
Expected-value rule                 & $0.728\pm0.014$ & --              & $0.664\pm0.010$ \\
Logistic regression (raw 4)         & $0.723\pm0.011$ & $0.619\pm0.010$ & $0.663\pm0.011$ \\
LDA (raw 4)                         & $0.718\pm0.012$ & $0.620\pm0.009$ & $0.656\pm0.013$ \\
Gaussian Naive Bayes (raw 4)        & $0.629\pm0.016$ & $0.688\pm0.017$ & $0.588\pm0.016$ \\
Linear SVM (raw 4)                  & $0.720\pm0.011$ & --              & --              \\
\midrule
LR + EV features                    & $0.785\pm0.015$ & $0.559\pm0.015$ & $0.704\pm0.014$ \\
\quad + domain and frame            & $0.814\pm0.014$ & $0.529\pm0.017$ & $0.732\pm0.016$ \\
\quad + person/session traits       & $0.812\pm0.014$ & $0.531\pm0.017$ & $0.732\pm0.012$ \\
\quad + frame$\times$EV interactions (no traits) & $0.816\pm0.011$ & $0.525\pm0.015$ & $0.734\pm0.012$ \\
\textbf{\quad + prospect-theory, trait interactions (final LR)} & $\mathbf{0.818\pm0.013}$ & $\mathbf{0.520\pm0.016}$ & $\mathbf{0.734\pm0.015}$ \\
ExtraTrees                          & $0.819\pm0.010$ & $0.519\pm0.013$ & $0.734\pm0.013$ \\
\textbf{Rank average LR + ExtraTrees (submitted)} & $\mathbf{0.820\pm0.011}$ & -- & -- \\
\bottomrule
\end{tabularx}
\end{table}

\begin{figure}[ht]
\centering
\includegraphics[width=\linewidth]{fig_overview.png}
\caption{(a) Share of risky choices per domain/frame cell in the training data. (b) ROC curves from pooled
out-of-fold predictions (pooled AUCs in the legend differ slightly from the fold means in Table~\ref{tab:main}).
(c) Calibration of the final logistic regression, 10 bins.}\label{fig:overview}
\end{figure}

\textbf{Experiment 1.} Raw-number models are no better than the EV rule ($\approx0.72$ vs.\ $0.728$), because EV is a
product of probability and outcome that a linear model on raw columns cannot form: the model class matters much
less than the features. Adding EV features raises AUC by $0.062$, and adding domain/frame by a further
$0.029\pm0.003$ (paired over folds; positive in all 5). Person and session traits do \emph{not} help
($-0.002\pm0.003$), and frame$\times$EV interactions and prospect-theory features add only $\approx0.002$ each, within
fold noise ($\mathrm{s.d.}\approx0.003$). The final logistic regression beats the EV rule by $0.090\pm0.011$.
Hypothesis 1 is supported.

\textbf{Experiment 2.} Random and participant-grouped splits give the same AUC for the final model ($0.818$ in
both; paired difference $-0.0004\pm0.022$): the random split is \emph{not} optimistic here, so that part of my
expectation was wrong. This is plausible because the model has no way to identify a person. Once a one-hot
\texttt{participant\_id} is added, the random split gives $0.839\pm0.013$, whereas the participant split stays at $0.818$
(no gain, as the test participants are unknown to the model). The paired gap is $+0.021\pm0.022$ over 5 folds, so it is
consistent in direction with the expectation but noisy given only five folds.

\textbf{Experiment 3.} Adding \texttt{response\_time\_ms} changes AUC by $-0.0001\pm0.0008$ (mean response time is 807~ms
for safe and 823~ms for risky choices), so excluding it loses nothing.

\textbf{Stronger models.} ExtraTrees alone is about equal to the logistic regression ($+0.0003\pm0.0045$), and the rank
average improves on the logistic regression by $+0.0018\pm0.0021$ (positive in 4 of 5 folds). This is about two standard errors
of the paired difference at best, so I regard the blend as at most a small improvement and use it because it was never
clearly worse. The submitted blend has validation AUC $0.820\pm0.011$.

\section{Error analysis and limitations}
\textbf{Where the model fails.} AUC differs strongly by cell (Table~\ref{tab:cell} in the Appendix): the easiest cell is
mixed gambles (blend AUC $0.90$), the hardest is gain problems in the gain frame ($0.68$). The model beats the EV rule in every
cell, with the largest gain in loss problems ($0.79$ vs.\ $0.73$). Performance is better in paid sessions (AUC $0.85$,
45\% of data) than hypothetical sessions ($0.79$). Per-participant accuracy of the final logistic regression ranges from $0.43$
to $0.95$ (mean $0.73$, s.d.\ $0.10$; Figure~\ref{fig:pacc} in the Appendix): the model has no information on whether a
new person is a strong risk taker or very careful, and traits such as the self-reported risk score did not
help. Calibration of the logistic regression is good (mean absolute gap between predicted and observed
share $0.007$, 10 bins), but the submitted blend is a rank score and is not calibrated.

\textbf{Limitations.} (i) The data are simulated, with random noise in the choice, so a perfect score is impossible and
nothing here says how real people behave; the findings only describe this generator. (ii) With 210 participants and
5 folds, differences below about $0.003$ AUC are not reliable; the standard deviations are spread over folds, not over
repeated experiments. (iii) The model choice and $C$ were selected on the same cross-validation, so $0.82$ may be
slightly optimistic for the test set. (iv) The frame effect (risky share $0.56$ vs.\ $0.40$ in the gain domain) is an
association in the data: I have not verified that the problems shown in the two frames are matched, and the frame is
only varied within the gain domain, so frame and domain cannot be separated outside it. (v) The prospect-theory features
use a few fixed parameter sets and no parameter is estimated, so they are only a rough approximation of the simulator.

\section{Conclusion}
For new participants, domain/frame and EV-based features clearly help: validation AUC rises from $0.72$ for the raw numbers
(and $0.73$ for the EV rule) to about $0.82$, whereas age, numeracy and self-reported risk taking add nothing for unseen
people. The result is limited to simulated data; a blend of logistic regression and ExtraTrees gives the best, but only marginally
better, score ($0.820\pm0.011$).

\section*{Reproducibility statement}
Code: The submitted ZIP file contains the notebook \texttt{Kanchan\_Bharti\_IML\_Mid\_term\_Project.ipynb}, a \texttt{README.md}, and a pinned \texttt{requirements.txt}. The repository is available at \href{https://github.com/Kanchan039/IML-2026-Track-B-Risky-Choice}{\texttt{IML-Mid\_Term\_Project}}.
Environment: Python 3.12 (pandas 2.3.3, numpy 2.2.6, scikit-learn 1.7.2, scipy 1.15.3, matplotlib 3.10.7).
Command: set \texttt{folder} in the first code cell and run all cells of the notebook (or \texttt{jupyter nbconvert --to notebook --execute}); it regenerates every number, table and figure in this report and saves \texttt{experimental\_results.csv} and \texttt{submission.csv} into that folder.
Seeds: folds and ExtraTrees use seed 0 and the output is identical across reruns. Runtime: a few minutes on one CPU core.
No neural networks or pretrained models were used. Generative AI (Claude) was used and is declared in the README.

\bibliographystyle{unsrtnat}
\begin{thebibliography}{6}
\bibitem{kahneman1979}
Kahneman, D. and Tversky, A. Prospect theory: an analysis of decision under risk.
\emph{Econometrica}, 47(2):263--291, 1979. doi: 10.2307/1914185.
\bibitem{tversky1981}
Tversky, A. and Kahneman, D. The framing of decisions and the psychology of choice.
\emph{Science}, 211(4481):453--458, 1981. doi: 10.1126/science.7455683.
\bibitem{tversky1992}
Tversky, A. and Kahneman, D. Advances in prospect theory: cumulative representation of uncertainty.
\emph{Journal of Risk and Uncertainty}, 5(4):297--323, 1992. doi: 10.1007/BF00122574.
\bibitem{roberts2017}
Roberts, D. R. et al. Cross-validation strategies for data with temporal, spatial, hierarchical, or phylogenetic structure.
\emph{Ecography}, 40(8):913--929, 2017. doi: 10.1111/ecog.02881.
\bibitem{geurts2006}
Geurts, P., Ernst, D. and Wehenkel, L. Extremely randomized trees.
\emph{Machine Learning}, 63(1):3--42, 2006. doi: 10.1007/s10994-006-6226-1.
\bibitem{pedregosa2011}
Pedregosa, F. et al. Scikit-learn: machine learning in Python.
\emph{Journal of Machine Learning Research}, 12:2825--2830, 2011.
\bibitem{fawcett2006}
Fawcett, T. An introduction to ROC analysis.
\emph{Pattern Recognition Letters}, 27(8):861--874, 2006. doi: 10.1016/j.patrec.2005.10.010.
\end{thebibliography}

\appendix
\section{Additional results}

\begin{table}[ht]
\centering
\caption{Experiments 2 and 3: effect of the split and of response time (final logistic regression, 5 folds, mean $\pm$ s.d.).}\label{tab:exp23}
\small
\begin{tabularx}{\linewidth}{@{}Xccc@{}}
\toprule
Setting & ROC-AUC & Log-loss & Accuracy \\
\midrule
Participant split (final LR)          & $0.818\pm0.013$ & $0.520\pm0.016$ & $0.734\pm0.015$ \\
Random split (final LR)               & $0.818\pm0.014$ & $0.519\pm0.015$ & $0.734\pm0.015$ \\
+ \texttt{participant\_id}, participant split & $0.818\pm0.013$ & $0.520\pm0.017$ & $0.734\pm0.014$ \\
+ \texttt{participant\_id}, random split      & $0.839\pm0.013$ & $0.495\pm0.015$ & $0.752\pm0.017$ \\
+ \texttt{response\_time\_ms}, participant split & $0.818\pm0.012$ & $0.520\pm0.016$ & $0.733\pm0.015$ \\
\bottomrule
\end{tabularx}
\end{table}

\begin{table}[ht]
\centering
\caption{Regularisation of the final logistic regression (participant split, 5 folds).}\label{tab:appC}
\small
\begin{tabular}{@{}lcccc@{}}
\toprule
$C$ & 0.003 & 0.03 & 0.1 (final) & 1 \\
\midrule
ROC-AUC & $0.8175\pm0.0133$ & $0.8182\pm0.0126$ & $0.8182\pm0.0127$ & $0.8180\pm0.0130$ \\
\bottomrule
\end{tabular}
\end{table}

\begin{table}[ht]
\centering
\caption{Paired AUC differences (same folds): mean $\pm$ s.d.\ over the 5 fold differences.}\label{tab:paired}
\small
\begin{tabularx}{\linewidth}{@{}Xc@{}}
\toprule
Comparison (A $-$ B) & $\Delta$AUC \\
\midrule
Final LR $-$ EV rule                        & $+0.090\pm0.011$ \\
Final LR $-$ LR on raw numbers              & $+0.095\pm0.006$ \\
+ domain/frame $-$ + EV features only       & $+0.029\pm0.003$ \\
+ person/session traits $-$ + domain/frame  & $-0.002\pm0.003$ \\
Final LR $-$ LR with frame$\times$EV interactions only & $+0.002\pm0.003$ \\
ExtraTrees $-$ final LR                     & $+0.0003\pm0.0045$ \\
Blend $-$ final LR                          & $+0.0018\pm0.0021$ \\
Random split $-$ participant split (final LR) & $-0.0004\pm0.022$ \\
Random $-$ participant split, with \texttt{participant\_id} & $+0.021\pm0.022$ \\
Final LR with response time $-$ final LR    & $-0.0001\pm0.0008$ \\
\bottomrule
\end{tabularx}
\end{table}

\begin{table}[ht]
\centering
\caption{AUC by domain/frame cell (pooled out-of-fold scores; the blend uses within-fold ranks).}\label{tab:cell}
\small
\begin{tabular}{@{}lrcccc@{}}
\toprule
Cell & $n$ & Share risky & EV rule & Final LR & Blend \\
\midrule
gain / gain frame  & 1\,731 & 0.398 & 0.632 & 0.672 & 0.679 \\
gain / loss frame  & 1\,649 & 0.557 & 0.722 & 0.766 & 0.764 \\
loss / none        & 2\,537 & 0.587 & 0.733 & 0.786 & 0.790 \\
mixed / none       & 2\,483 & 0.383 & 0.858 & 0.898 & 0.901 \\
\bottomrule
\end{tabular}
\end{table}

\begin{figure}[ht]
\centering
\includegraphics[width=0.6\linewidth]{fig_participant_acc.png}
\caption{Out-of-fold accuracy per participant for the final logistic regression (210 participants).}\label{fig:pacc}
\end{figure}

\end{document}
